\subsection{Data Sources: Actual and Assumed Data}\label{sec:data}

\begin{table}[H]
\centering
\caption{Sources of the data, separating actual data from assumptions.}\label{tab:sources}
\vspace{4pt}
\small
\begin{tabularx}{\textwidth}{l l >{\raggedright\arraybackslash}X}
\toprule
\textbf{Item} & \textbf{Type} & \textbf{Source / justification} \\
\midrule
Locations (V0--V9) & Actual & Real places in Cagayan de Oro City, located on Google
  Maps~\cite{googlemaps} and OpenStreetMap~\cite{osm}. \groupnote{Confirm each pin in Google Maps
  and add screenshots.} \\
Coordinates & Actual & OpenStreetMap, retrieved with the Nominatim geocoder~\cite{osm}. \\
Road distances $d(u\to v)$ & Actual & Shortest driving route on the OpenStreetMap road network,
  computed by the OSRM routing engine~\cite{luxen2011osrm} (\texttt{code/tools/build\_dataset.py}).
  \groupnote{Spot-check several pairs against Google Maps driving distances.} \\
Symmetric weights $w(u,v)$ & Assumed & Average of the two directions (Section~\ref{sec:edges}). \\
Travel mode, no traffic or time windows & Assumed & Simplifications listed in Section~\ref{sec:problem}. \\
Choice of the nine stops and the start & Assumed & Scenario chosen by the group. \\
\bottomrule
\end{tabularx}
\end{table}

All the data are in plain CSV files under \texttt{code/data/}:
\begin{itemize}
  \item \texttt{locations.csv}: the vertices (Table~\ref{tab:vertices});
  \item \texttt{distance\_matrix\_km.csv}: the edge weights (Table~\ref{tab:matrix});
  \item \texttt{raw\_directed\_km.csv}: the raw directed distances, before averaging.
\end{itemize}
The program reads these files directly, and every table
and figure in this report is generated from the same run of the program, so the graph, the data, the
program input and the reported results are consistent.
